% GENERATED FILE. Edit canonical lessons, then run npm run generate:books.
% canonical-source-hash: fnv1a64:140392c9d012986a
% canonical-lessons: SA-C156-priyah, SA-C156-apriyah, SA-C156-santah, SA-C156-candah, SA-C156-krurah

\chapter{Dear, Unpleasant, Calm, Fierce, Cruel}
\label{ch:sa-dear-unpleasant-calm-fierce-cruel}
\hlchaptermodality{156}

\begin{chapteropening}
\textbf{By the end of this chapter:} \emph{I can say dear, unpleasant, calm, fierce and cruel.}

Complete the last lesson of chapter 156: I can say dear, unpleasant, calm, fierce and cruel.
\end{chapteropening}

\section[\texorpdfstring{\sk{प्रियः} (priyaḥ)}{प्रियः (priyaḥ)}]{\sk{प्रियः} (priyaḥ) — dear, beloved}

\label{lesson:SA-C156-priyah}



\begin{quote}
Before the new one: say the Sanskrit for chilly, then the Sanskrit for harsh.
\end{quote}

\subsection*{You'll want to know: \sk{प्रियः}}
\textbf{\sk{प्रियः}} — \emph{priyaḥ} — \textquotedblleft{}dear, beloved\textquotedblright{}.

\begin{grammarlens}[title={What you've built}]
One more word for this chapter.
\end{grammarlens}

\subsection*{Your turn}
\begin{itemize}
\raggedright
  \item \emph{Say it:} \emph{priyaḥ}
  \item \emph{Say it:} \emph{priyaḥ}, once more
  \item \emph{Say it:} say \emph{priyaḥ}
  \item \emph{Recall:} say the Sanskrit for chilly, then the Sanskrit for harsh, then say \emph{priyaḥ} again
\end{itemize}

\begin{tcolorbox}[breakable,colback=teal!4,colframe=teal!35!black,title={Before you move on}]
What does \sk{प्रियः} mean? (\textquotedblleft{}Dear, beloved\textquotedblright{}.) Say it once more.
\end{tcolorbox}
\section[\texorpdfstring{\sk{अप्रियः} (apriyaḥ)}{अप्रियः (apriyaḥ)}]{\sk{अप्रियः} (apriyaḥ) — unpleasant}

\label{lesson:SA-C156-apriyah}



\begin{quote}
Before the new one: say the Sanskrit for harsh, then the Sanskrit for dear.
\end{quote}

\subsection*{You'll want to know: \sk{अप्रियः}}
\textbf{\sk{अप्रियः}} — \emph{apriyaḥ} — \textquotedblleft{}unpleasant\textquotedblright{}.

\begin{grammarlens}[title={What you've built}]
One more word for this chapter.
\end{grammarlens}

\subsection*{Your turn}
\begin{itemize}
\raggedright
  \item \emph{Say it:} \emph{apriyaḥ}
  \item \emph{Say it:} \emph{apriyaḥ}, once more
  \item \emph{Say it:} say \emph{apriyaḥ}
  \item \emph{Recall:} say the Sanskrit for harsh, then the Sanskrit for dear, then say \emph{apriyaḥ} again
\end{itemize}

\begin{tcolorbox}[breakable,colback=teal!4,colframe=teal!35!black,title={Before you move on}]
What does \sk{अप्रियः} mean? (\textquotedblleft{}Unpleasant\textquotedblright{}.) Say it once more.
\end{tcolorbox}
\section[\texorpdfstring{\sk{शान्तः} (śāntaḥ)}{शान्तः (śāntaḥ)}]{\sk{शान्तः} (śāntaḥ) — calm, quiet}

\label{lesson:SA-C156-santah}



\begin{quote}
Before the new one: say the Sanskrit for dear, then the Sanskrit for unpleasant.
\end{quote}

\subsection*{You'll want to know: \sk{शान्तः}}
\textbf{\sk{शान्तः}} — \emph{śāntaḥ} — \textquotedblleft{}calm, quiet\textquotedblright{}.

\begin{grammarlens}[title={What you've built}]
One more word for this chapter.
\end{grammarlens}

\subsection*{Your turn}
\begin{itemize}
\raggedright
  \item \emph{Say it:} \emph{śāntaḥ}
  \item \emph{Say it:} \emph{śāntaḥ}, once more
  \item \emph{Say it:} say \emph{śāntaḥ}
  \item \emph{Recall:} say the Sanskrit for dear, then the Sanskrit for unpleasant, then say \emph{śāntaḥ} again
\end{itemize}

\begin{tcolorbox}[breakable,colback=teal!4,colframe=teal!35!black,title={Before you move on}]
What does \sk{शान्तः} mean? (\textquotedblleft{}Calm, quiet\textquotedblright{}.) Say it once more.
\end{tcolorbox}
\section[\texorpdfstring{\sk{चण्डः} (caṇḍaḥ)}{चण्डः (caṇḍaḥ)}]{\sk{चण्डः} (caṇḍaḥ) — fierce}

\label{lesson:SA-C156-candah}



\begin{quote}
Before the new one: say the Sanskrit for unpleasant, then the Sanskrit for calm.
\end{quote}

\subsection*{You'll want to know: \sk{चण्डः}}
\textbf{\sk{चण्डः}} — \emph{caṇḍaḥ} — \textquotedblleft{}fierce\textquotedblright{}.

\begin{grammarlens}[title={What you've built}]
One more word for this chapter.
\end{grammarlens}

\subsection*{Your turn}
\begin{itemize}
\raggedright
  \item \emph{Say it:} \emph{caṇḍaḥ}
  \item \emph{Say it:} \emph{caṇḍaḥ}, once more
  \item \emph{Say it:} say \emph{caṇḍaḥ}
  \item \emph{Recall:} say the Sanskrit for unpleasant, then the Sanskrit for calm, then say \emph{caṇḍaḥ} again
\end{itemize}

\begin{tcolorbox}[breakable,colback=teal!4,colframe=teal!35!black,title={Before you move on}]
What does \sk{चण्डः} mean? (\textquotedblleft{}Fierce\textquotedblright{}.) Say it once more.
\end{tcolorbox}
\section[\texorpdfstring{\sk{क्रूरः} (krūraḥ)}{क्रूरः (krūraḥ)}]{\sk{क्रूरः} (krūraḥ) — cruel}

\label{lesson:SA-C156-krurah}



\begin{quote}
Before the new one: say the Sanskrit for calm, then the Sanskrit for fierce.
\end{quote}

\subsection*{You'll want to know: \sk{क्रूरः}}
\textbf{\sk{क्रूरः}} — \emph{krūraḥ} — \textquotedblleft{}cruel\textquotedblright{}.

\begin{grammarlens}[title={What you've built}]
One more word for this chapter.
\end{grammarlens}

\subsection*{Your turn}
\begin{itemize}
\raggedright
  \item \emph{Say it:} \emph{krūraḥ}
  \item \emph{Say it:} \emph{krūraḥ}, once more
  \item \emph{Say it:} say \emph{krūraḥ}
  \item \emph{Recall:} say the Sanskrit for calm, then the Sanskrit for fierce, then say \emph{krūraḥ} again
\end{itemize}

\begin{tcolorbox}[breakable,colback=teal!4,colframe=teal!35!black,title={Before you move on}]
What does \sk{क्रूरः} mean? (\textquotedblleft{}Cruel\textquotedblright{}.) Say it once more.
\end{tcolorbox}
